\chapter{Theory}

\section{First Section}
Some theory... 
\begin{equation}
	a^2 + b^2 = c^2
\end{equation}

\section{Second Section}
Lorem ipsum

\subsection{A Subsection}
\label{subsec: a-subsection}
\textit{Text kurzivou. }
\textbf{Text boldom. }
\emph{Vyznamovo zvyrazeny text } \newline
\section{}
Lorem ipsum 

\subsection{A Subsection with Figure}


Lorem ipsum 
leo sit amet mattis ullamcorper.

\begin{figure}[b]
	\centering
	\includegraphics[width = 0.5\textwidth]{images/fchpt_logo_color}
	\caption{A figure caption.}
	\label{fig:label}
\end{figure}

\begin{tikzpicture}
    \draw[-{Latex}] (0,0) -- (4,0) node[right] {$x$};
    \draw[-{Latex}] (0,0) -- (0,3) node[above] {$y$};
    \draw[red, very thick] (0,0) -- (1,1.4) -- (2.1,0.9) -- (2.8,2.2);
    \node[draw, rounded corners, fill=cyan] at (2.8,2.5) {merane udaje};
\end{tikzpicture}

\begin{equation*}
    f(x)=
\begin{cases}
    x^2 & \text{ak } x> 0, \\
    -x, & \text{ak } x<0, \\
    0, & \text{ak } x=0
\end{cases}
\end{equation*}


Maecenas posuere 
\begin{table}[]
\centering
\caption{Zoznam studentov a predmetov FCHPT}
\label{tab:stunt1-uiam}
\begin{tabular}{lllll}
 & 6     &  &  &  \\
 & twkjv &  &  &  \\
 &       &  &  &  \\
 &       &  &  & 
\end{tabular}
\end{table}
